// Emboss — directional emboss via offset sampling
//
// Creates a raised/carved appearance by computing the luminance difference
// between the current pixel and an offset pixel. The difference is added
// to mid-gray (0.5), so raised edges appear bright and recessed edges
// appear dark — like light hitting a textured surface from one direction.
//
// Connect @image to any RGB input.

f$angle = 45.0;         // light direction in degrees (0 = right, 90 = down)
f$strength = 1.0;       // emboss intensity multiplier
f$distance = 1.0;       // offset distance in pixels
f$blend = 0.0;          // blend with original (0 = pure emboss, 1 = overlay)

vec3 col = @image;

// ── Compute offset direction ────────────────────────────────────────
float rad = $angle * PI / 180.0;
float dx = cos(rad) * $distance;
float dy = sin(rad) * $distance;

// Convert pixel offset to UV offset
float du = dx / iw;
float dv = dy / ih;

// ── Sample at offset position ───────────────────────────────────────
vec3 offset_col = sample(@image, u + du, v + dv);

// ── Compute luminance difference ────────────────────────────────────
float lum_center = luma(col);
float lum_offset = luma(offset_col);
float diff = (lum_center - lum_offset) * $strength;

// ── Build emboss result on mid-gray base ────────────────────────────
// Mid-gray (0.5) is neutral; positive diff = highlight, negative = shadow
vec3 emboss = vec3(0.5 + diff);

// ── Optional blend with original color ──────────────────────────────
// When blend > 0, the emboss effect is mixed over the original image,
// adding a textured/relief look while preserving color.
vec3 result = lerp(emboss, col + vec3(diff), $blend);

@OUT = clamp(result, 0.0, 1.0);
